% !TEX root = ../main_ext_1.tex

% INTRODUCTION PATRICK
\section{Introduction}
Impulse responses are primarily estimated via vector autoregressive (VAR) or local projection (LP) models. While LPs and VARs estimate the same population impulse responses under general conditions \parencite{plagborg2021local, fusari2026more}, their finite-sample behavior differs. At equal lag lengths, LPs are viewed as less sensitive to misspecification \parencite{jorda2005estimation, ramey2016macroeconomic}---a property formalised by \textcite{olea2024double}---whereas VARs appear more statistically efficient \parencite{li2024local, olea2025local}. \textcite{ludwig2026local} analytically explains this bias-variance tradeoff: an LP mathematically represents a sequence of VARs of increasing lag length. He thus argues for comparing estimators by effective complexity rather than equal lag length.

Our paper asks two questions. First, where does a fixed $\text{LP}(4)$ sit relative to the complexity-adjusted VAR frontier under correct specification, and does the conventional equal-lag ranking survive? We benchmark a fixed $\text{LP}(4)$ against progressively higher-order VAR specifications \parencite{de2025long} across empirically relevant sample sizes \parencite[see e.g.][]{herbst2024bias}. Second, once controlled dynamic misspecification is introduced, does any genuine LP advantage emerge in point estimation, inference, or both? By augmenting the baseline with a moving-average component to yield a $\text{VARMA}(4,1)$ DGP, we track how $\text{LP}(4)$ and the VAR sweep respond. This is the setting where LPs are conventionally expected to dominate, as their conjectured robustness should bind precisely where finite-order VARs are misspecified \parencite{jorda2005estimation, li2024local,olea2024double}.

Under correct specification, the point-estimation advantage of equal-lag VARs persists, but largely as an artifact of mismatched model complexity. A fixed $\text{LP}(4)$ operates like a highly parameterized system, sitting high on the VAR frontier (beyond $q \approx 15$). This low-order efficiency comes at a cost in inference, as $\text{VAR}(4)$ coverage deteriorates sharply near the unit root while LPs hold near-nominal coverage. Introducing dynamic misspecification does not overturn the point-estimation ranking; the equal-lag $\text{VAR}(4)$ still strictly dominates $\text{LP}(4)$ in root mean squared error. It instead sharpens the inferential divergence, where LPs retain their advantage. However, this disparity stems strictly from the lag restriction. When evaluated against appropriately higher-order VARs on the complexity frontier, the finite-sample gap systematically closes. A sufficiently flexible VAR successfully absorbs the unmodeled moving-average dynamics, catching up to the $\text{LP}(4)$ in bias reduction and matching its robust coverage levels, even near the non-stationary boundary.

The remainder of the paper is structured as follows. Section 2 reviews vector autoregressions 
and local projections, the empirical literature on their bias-variance trade-off at equal 
lag lengths, and their asymptotic and finite-sample equivalence. Section 3 presents the 
baseline simulation study under correct specification, Section 4 introduces dynamic 
misspecification, and Section 5 concludes.

\FloatBarrier